\clearpage
\VerdierPhysicalPageAnchor{0113}%
\begingroup
\small
\setlength{\abovedisplayskip}{4pt plus 2pt minus 2pt}%
\setlength{\belowdisplayskip}{4pt plus 2pt minus 2pt}%

\VerdierSemanticAnchor{ch2-1-1-5-additivity-announcement}
\noindent (1.2.7) において、コホモロジー的関手が加法的であり、多重完全関手が
多重加法的であることを示す。

\VerdierSemanticAnchor{ch2-1-1-5-cohomology-notation}
$H\colon\mathcal D\longrightarrow\mathcal A$ をコホモロジー的関手とする。
通常、$H^0(X)=H(X)$、$H^n(X)=H(X[n])$ と書く。公理 TRII を用いると、
コホモロジー的関手は任意の完全三角

\VerdierSemanticAnchor{ch2-1-1-5-long-exact-setup}
\[
  X\xrightarrow{u}Y\xrightarrow{v}Z\xrightarrow{w}X[1]
\]
に長完全列

\VerdierSemanticAnchor{ch2-1-1-5-1}
\begin{equation}
  \cdots\longrightarrow H^{-1}(Z)
  \xrightarrow{H^{-1}(-w)}H^0(X)
  \xrightarrow{H^0(u)}H^0(Y)
  \xrightarrow{H^0(v)}H^0(Z)
  \xrightarrow{H^0(w)}\cdots
  \tag{1.1.5.1}
\end{equation}
を対応させることが分かる。

\VerdierSemanticAnchor{ch2-1-1-6}
\medskip
\noindent\textbf{1.1.6.}\quad
$\mathcal D$ を三角圏とする。通常、次の用語を用いる。$\mathcal D$ の三角

\VerdierSemanticAnchor{ch2-1-1-6-antidistinguished-triangle}
\[
  X\xrightarrow{u}Y\xrightarrow{v}Z\xrightarrow{w}X[1]
\]
が\emph{反完全}であるとは、三角
\[
  X\xrightarrow{u}Y\xrightarrow{v}Z\xrightarrow{-w}X[1]
\]
が完全三角であることをいう。

\VerdierSemanticAnchor{ch2-1-1-6-antiexact-functor}
三角圏 $\mathcal D'$ から $\mathcal D$ への $\mathbb Z$-関手が
\emph{反完全}であるとは、それが $\mathcal D'$ の完全三角を
$\mathcal D$ の反完全三角へ送ることをいう。

\VerdierSemanticAnchor{ch2-1-1-6-exact-functor-notation}
通常、$\mathcal D'$ から $\mathcal D$ への完全関手を記す際には、シフトと
可換する同型を記号から省略する。

\VerdierSemanticAnchor{ch2-1-1-7}
\medskip
\noindent\textbf{1.1.7.}\quad
$\mathcal C$ を厳密な $\mathbb Z$-圏（第 I 章、1.2.2）とし、

\VerdierSemanticAnchor{ch2-1-1-7-translations}
\[
  X\longmapsto X[n],\qquad n\in\mathbb Z,
\]
を $\mathcal C$ のシフト関手とする。$\mathcal C^\circ$ を反対
$\mathbb Z$-圏（第 I 章、1.3.9）とする。

\VerdierSemanticAnchor{ch2-1-1-7-opposite-morphism}
$\mathcal C$ の任意の射 $u$ に対し、$\mathcal C^\circ$ の対応する射を
$((u))$ と書く。$\mathcal C^\circ$ の対象の集合と射の集合は、それぞれ
$\mathcal C$ の対象の集合と射の集合に等しいので、この記法は厳密には
濫用である。$\mathcal C^\circ$ のシフト関手を

\VerdierSemanticAnchor{ch2-1-1-7-opposite-translations}
\[
  X\longmapsto X[[n]],\qquad n\in\mathbb Z,
\]
と書く。

\VerdierSemanticAnchor{ch2-1-1-7-translation-identity}
$\mathcal C$ の任意の対象 $X$ と任意の整数 $n\in\mathbb Z$ に対し
（第 I 章、1.3.9）、
\[
  X[-n]=X[[n]]
\]
であったことを思い出そう。

\VerdierSemanticAnchor{ch2-1-1-7-opposite-category}
三角圏 $\mathcal D$ の\emph{反対三角圏}とは、圏 $\mathcal D$ の厳密な
反対 $\mathbb Z$-圏 $\mathcal D^\circ$（第 I 章、1.3.9）に、次の完全三角の
族を入れたものと定義する。$\mathcal D^\circ$ の三角

\VerdierSemanticAnchor{ch2-1-1-7-opposite-triangle}
\[
  X[[-1]]\xrightarrow{((w))}Z\xrightarrow{((v))}Y
  \xrightarrow{((u))}X
\]
が完全三角であるための必要十分条件は、圏 $\mathcal D$ の対応する三角

\VerdierSemanticAnchor{ch2-1-1-7-corresponding-antitriangle}
\[
  X\xrightarrow{u}Y\xrightarrow{v}Z\xrightarrow{w}X[1]
\]
が反完全三角であることである。

\VerdierSemanticAnchor{ch2-1-1-7-triangulated-axioms}
このように定めた $\mathcal D^\circ$ の三角の族が三角圏の公理を満たすことは、
容易に示される。
\endgroup
